\documentclass[11pt]{article}
\usepackage[margin=1in]{geometry}
\usepackage{amsmath,amssymb,amsthm,booktabs}
\usepackage{hyperref}
\hypersetup{hidelinks,pdftitle={Four-zero tail-palindromic polynomials over F2},pdfauthor={Baohua Li}}
\newtheorem{proposition}{Proposition}[section]
\newtheorem{lemma}[proposition]{Lemma}
\newtheorem{theorem}[proposition]{Theorem}
\newtheorem{corollary}[proposition]{Corollary}
\theoremstyle{definition}
\newtheorem{conjecture}[proposition]{Conjecture}
\theoremstyle{remark}
\newtheorem{remark}[proposition]{Remark}
\title{\Large Four-zero tail-palindromic polynomials over $\mathbb F_2$:\\
reciprocal obstructions and quadratic-transform intersections}
\author{Baohua Li}
\date{29 October 2026}
\begin{document}
\maketitle
\begin{abstract}
We study monic even-degree polynomials over $\mathbb F_2$ whose nonleading coefficients form a palindrome with constant term one and exactly four zeros. We completely classify their intersections with full-degree quadratic rational transformations. The translated $Q$ branch consists of two infinite shape families and three degree-14 exceptions; the translation-invariant branch consists of four reducible polynomials at degrees 14 and 30. More generally, a translation-invariant tail palindrome of even degree $n\ge6$ has $n+2$ a power of two and is reducible whenever its number of missing tail coefficients is divisible by four. The classification determines the boundary of quadratic degree doubling: every four-zero output has degree $2$ modulo $4$, so iterated doubling from an even input degree cannot produce a member of the family. We also determine the exact reciprocal common divisor and count the position pairs surviving this obstruction. Reproducible arithmetic computations identify 2663 primitive polynomials at the 98 even degrees from 6 to 200 and verify the first solutions through degree 300. Conjecture $C_4$ on primitive existence at every even degree remains open.
\end{abstract}

\section{The family and the question}\label{sec:intro}
Quadratic rational transformations are a standard tool for constructing
irreducible polynomials of doubled degree \cite{MP17}. This paper asks
which outputs can also satisfy a rigid coefficient constraint: a
palindromic tail with exactly four zero coefficients.

Put $n=2k\ge6$ and $S_n(x)=1+x+\cdots+x^n$. We consider
\begin{equation}\label{eq:candidate}
 f_{n,i,j}(x)=S_n(x)+x^i+x^j+x^{n-1-i}+x^{n-1-j},
 \qquad 1\le i<j\le k-1,
\end{equation}
over $\mathbb F_2$. These are exactly the monic degree-$n$ polynomials
with constant term one whose nonleading coefficient vector is a
palindrome with four zero entries. The restrictions on $i,j$ exclude
the constant position. There are $\binom{k-1}{2}$ candidates.
Let $N_4(n)$ denote the number of primitive polynomials in this family.
Here primitive means irreducible with a root of multiplicative order
$2^n-1$; standard background is in \cite{LN}.

\begin{conjecture}[$C_4$]\label{conj:C4}
For every even $n\ge6$, one has $N_4(n)>0$.
\end{conjecture}

Primitive polynomials with prescribed coefficients are studied systematically
by Cohen \cite{Coh06}. In the present family the field is fixed, all but four
tail coefficients are one, and the four missing positions are coupled by
reversal. The number of coefficient constraints therefore grows with the
degree. A tail palindrome can also be encoded by the self-reciprocal
polynomial $xf(x)+1$ of odd degree $n+1$. This describes its support;
irreducibility and primitivity refer throughout to $f$ itself.

The main results completely classify the full-degree quadratic outputs.
Theorems~\ref{thm:quadratic} and~\ref{thm:translation-classification}
give the translated $Q$ and translation-invariant lists. A trace obstruction
also yields uniform reducibility when the number of missing tail
coefficients is divisible by four
(Proposition~\ref{prop:translation-zero-obstruction}).
Corollary~\ref{cor:all-quadratic} combines the branches and rules out
iterated doubling from an even input degree. The proof converts the
coefficient constraints into coupled translation identities, whose solutions
are determined by binary submask arguments.

Section~\ref{sec:reciprocal-screen} gives a complementary arithmetic
screen: an exact reciprocal common divisor and a closed count of the
surviving position pairs (Propositions~\ref{prop:reciprocal-gcd}
and~\ref{prop:reciprocal-survivors}).

Section~\ref{sec:computations} gives reproducible primitive counts and
first witnesses. Section~\ref{sec:limits} develops the questions suggested
by the classification, while Appendix~\ref{sec:classical-descriptions}
describes coefficient multiplicities and finite-field root parameters.

\section{The reciprocal obstruction}
\label{sec:reciprocal-screen}

Tail symmetry turns polynomial reciprocity into an explicit common-factor
test. In the four-zero family, the forced factor is determined by two
integer gcds, and the number of surviving position pairs has a closed
formula. The calculation uses the binomial gcd identity \cite{LN} and
unit counts expressed by the second Schemmel function \cite[Section~1]{Def20}.

\begin{lemma}[Reciprocity and multiplicities]
\label{lem:reciprocal-multiplicity}
Let \(f=x^n+\sum_{r=0}^{n-1}c_rx^r\), where \(n\) is even,
\(c_0=1\), and \(c_r=c_{n-1-r}\). Write
\(f^*(x)=x^nf(1/x)\). Then
\[
 f^*=xf+x^{n+1}+1,\qquad
 \gcd(f,f^*)=\gcd(f,x^{n+1}+1).
\]
This gcd is squarefree. In particular, every self-reciprocal
irreducible factor of \(f\) occurs with multiplicity one. If a
non-self-reciprocal irreducible \(g\) and \(g^*\) both divide \(f\),
the smaller of their multiplicities is one.
\end{lemma}
\begin{proof}
The coefficient symmetry gives the displayed identity, hence the gcd
identity by subtraction. The derivative of \(x^{n+1}+1\) is \(x^n\),
which is coprime to it; thus its divisors are squarefree. A
self-reciprocal factor has the same multiplicity in \(f\) and \(f^*\).
For a reciprocal pair the two multiplicities are interchanged, so
their minimum is the multiplicity in the gcd. This proves the claims.
\end{proof}

\begin{proposition}[Exact reciprocal common divisor]
\label{prop:reciprocal-gcd}
For \(f=f_{n,i,j}\) from \eqref{eq:candidate}, put
\[
 N=n+1,\quad \Delta=j-i,\quad R=n-1-i-j,\quad
 a=\gcd(N,\Delta),\quad b=\gcd(N,R),\quad h=\gcd(a,b).
\]
Then, as an exact polynomial quotient in \(\mathbb F_2[x]\),
\begin{equation}\label{eq:reciprocal-gcd}
 \gcd(f,f^*)=
 \frac{(x^a+1)(x^b+1)}{(x+1)(x^h+1)},\qquad
 \deg\gcd(f,f^*)=a+b-h-1.
\end{equation}
Consequently, irreducibility of \(f\) requires \(a=b=1\).
\end{proposition}
\begin{proof}
The nonleading terms of \(f\) occur in pairs, so \(f(1)=1\).
By Lemma~\ref{lem:reciprocal-multiplicity}, with
\(S_n=(x^N+1)/(x+1)\),
\[
 \gcd(f,f^*)=\gcd(f,S_n)=\gcd(S_n,P),\qquad
 P=x^i(1+x^\Delta)(1+x^R).
\]
The oddness of \(N\) makes \(S_n\) squarefree. For positive $u,v$,
the Euclidean algorithm on the exponents gives
$\gcd(x^u+1,x^v+1)=x^{\gcd(u,v)}+1$ over $\mathbb F_2$.
Applying this identity yields
\[
 \gcd(S_n,1+x^\Delta)=\frac{x^a+1}{x+1},\qquad
 \gcd(S_n,1+x^R)=\frac{x^b+1}{x+1}.
\]
Since \(S_n\) is squarefree, its gcd with their product is their
least common multiple. Their intersection is
\((x^h+1)/(x+1)\), proving \eqref{eq:reciprocal-gcd}.
If this gcd were all of \(f\), monicity and equal degrees would give
\(f=f^*\), hence \((x+1)f=x^N+1\) and \(f=S_n\), contrary to its
four zero coefficients. Thus a nonconstant gcd is a proper divisor.
The right side is constant exactly when \(a=b=1\).
\end{proof}

\begin{corollary}[An explicit infinite reducible family]
\label{cor:infinite-reducible}
For every odd \(m\ge3\), the valid position pair
\((n,i,j)=(3m-1,m-2,m)\) satisfies
\begin{equation}\label{eq:infinite-reducible}
\begin{split}
 f_{3m-1,m-2,m}(x)
 ={}&(1+x+\cdots+x^{m-1})\\
 &\quad\cdot(x^{2m}+x^{m+1}+x^{m-1}+x^{m-2}+1).
\end{split}
\end{equation}
In particular it is reducible, and the displayed forced divisor
has unbounded degree \(m-1\).
\end{corollary}
\begin{proof}
The inequalities \(1\le m-2<m\le(3m-3)/2\) make the positions valid.
Here \(N=3m\), \(\Delta=2\), and \(R=m\), so
\(a=1\), \(b=m\) in Proposition~\ref{prop:reciprocal-gcd}.
For the explicit quotient, group the geometric sum as
\(S_{3m-1}=S_{m-1}(1+x^m+x^{2m})\). The four deleted terms sum to
\[
 x^{m-2}(1+x^2)(1+x^m)=S_{m-1}x^{m-2}(x+1)^3.
\]
Adding these expressions and cancelling the two \(x^m\) terms
inside the second factor proves \eqref{eq:infinite-reducible}.
\end{proof}

\begin{proposition}[Number of positions surviving the screen]
\label{prop:reciprocal-survivors}
For odd \(N=n+1\), let \(\varphi\) be Euler's totient and put
\[
 \mathfrak S_2(N)=N\prod_{p\mid N}\left(1-\frac2p\right).
\]
The number of four-zero position pairs with \(\gcd(f,f^*)=1\) is
\begin{equation}\label{eq:screen-survivors}
 E_N=\binom{\varphi(N)/2}{2}-\frac{\mathfrak S_2(N)-1}{2}.
\end{equation}
Thus \(N_4(n)\le E_{n+1}\). This is only an upper bound;
\(E_{n+1}>0\) does not imply \(N_4(n)>0\).
\end{proposition}
\begin{proof}
Set \(u=i+1\), \(v=j+1\). The two coprimality conditions in
Proposition~\ref{prop:reciprocal-gcd} are precisely that \(u-v\)
and \(u+v\) be units modulo \(N\). Since \(2\) is invertible,
the map \((u,v)\mapsto(u-v,u+v)\) is a bijection on
\((\mathbb Z/N\mathbb Z)^2\), giving \(\varphi(N)^2\) ordered
pairs before excluding zero coordinates. The cases \(u=0\) and
\(v=0\) contribute \(\varphi(N)\) each and are disjoint.
The independent sign changes of \(u,v\) and their interchange
act freely on the remaining pairs: admissibility excludes
\(u=\pm v\), and a nonzero coordinate cannot equal its negative.
Dividing by eight gives
\(\varphi(N)(\varphi(N)-2)/8\) unordered pairs of distinct
nonzero sign classes.

Choose representatives \(1,\ldots,k\) of those classes. The
allowed positions use \(2,\ldots,k\), so the class represented
by \(1\) must be removed. The Chinese remainder theorem counts
\(\mathfrak S_2(N)\) residues \(v\) for which both \(v-1\) and
\(v+1\) are units: modulo each prime \(p\mid N\) two residues
are forbidden. One residue is \(v=0\); the others form pairs
\(\{v,-v\}\). Hence the removed class has
\((\mathfrak S_2(N)-1)/2\) neighbors. Subtraction proves
\eqref{eq:screen-survivors}.
\end{proof}

\begin{remark}[Factorization examples and fixed positions]
\label{rem:screen-scope}
For example,
\[
 f_{8,1,2}=(x^3+x^2+1)(x^5+x^2+1)
\]
has trivial reciprocal gcd and is reducible. Repeated factors also occur:
\[
 f_{14,2,5}=(x^2+x+1)(x^6+x^4+x^3+x+1)^2.
\]
The squarefreeness hypothesis in the classical
Stickelberger--Swan factor-parity criterion \cite{Swan62} must
therefore be checked. An odd number of factors still need not mean
irreducibility: \(f_{12,4,5}\) has three distinct factors of degrees
\(3,4,5\).

For fixed \(i,j\), one also has
\[
 (x+1)f_{n,i,j}=x^{n-1-j}A_{i,j}+B_{i,j},
\]
where \(A_{i,j}=x^{j+2}+(x+1)(x^{j-i}+1)\) and
\(B_{i,j}=1+(x+1)(x^i+x^j)\) are fixed polynomials.
This places fixed-position factor-parity questions in the
classical framework of Fredricksen, Hales and Sweet \cite{FHS86}.
\end{remark}


\section{Quadratic-transform intersections}
\label{sec:quadratic}
\subsection{The translated \texorpdfstring{$Q$}{Q} branch}
For a degree-$m$ polynomial $h$, define
\begin{equation}\label{eq:quadratic-transform}
 \mathcal T(h)(x)=(x+1)^m h\!\left(\frac{x^2}{x+1}\right),\qquad
 \mathcal M_n(f)(x)=(x+1)^n f\!\left(\frac{x}{x+1}\right).
\end{equation}
The map $\mathcal T$ is the classical $Q$-transform followed by the
translation $x\mapsto x+1$; explicitly, $\mathcal T(h)(x)=h^Q(x+1)$.
The equivalence between
$\mathcal M_{2m}(f)=f$ and $f=\mathcal T(h)$ is a special case of
Mattarei and Pizzato \cite[Lemma 9]{MP17}: take $a=b=1,c=0$,
numerator $x^2$ and denominator $x+1$ in their notation.
The nondegeneracy conditions hold, and the representation is unique;
indeed the summands $x^{2j}(x+1)^{m-j}$ have distinct degrees $m+j$.
More general invariant representations are given by Reis
\cite[Theorem 1.3]{Reis20}. We combine these invariant representations
with the four-zero coefficient condition to determine every compatible
output, including reducible polynomials.

Write $\Delta a=a(x+1)+a(x)$, and use reversal with respect to a
specified formal degree when the leading coefficient may vanish.

\begin{lemma}[A descent for tail palindromes]\label{lem:descent}
Let $f=x^{2k}+\sum_{j=0}^{2k-1}c_jx^j$, with $c_0=1$ and
$c_j=c_{2k-1-j}$. Set $d=k-1$ and
\[
 e=\sum_{a=0}^{k-1}c_{2a}x^a,\qquad o=x^d e(1/x).
\]
Then
\begin{equation}\label{eq:descent}
 \mathcal M_{2k}(f)=f
 \quad\Longleftrightarrow\quad
 \Delta e=0\ \hbox{and}\ e=\Delta(xo).
\end{equation}
If the tail of $f$ has $2s$ zeros, the formal coefficient vector of
$e$ has $s$ zeros.
\end{lemma}
\begin{proof}
Tail symmetry gives $f=x^{2k}+e^2+xo^2$. Put
$p=xf+x^{2k+1}=xe^2+x^2o^2$, and homogenize $p$ to formal degree
$2k+1$, obtaining $P(X,Y)$. The tail condition is exactly
$P(X,Y)=P(Y,X)$. Moreover, $\mathcal M_{2k}(f)=f$ is equivalent
to $P(X,X+Y)=P(X,Y)$. Conjugating this substitution by the
interchange of $X,Y$ gives $(X,Y)\mapsto(X+Y,Y)$. Hence, for
this symmetric $P$, the condition is equivalent to $\Delta p=0$.
Direct expansion gives
\[
 \Delta p=x(\Delta e)^2+
 \bigl(e(x+1)+\Delta(xo)\bigr)^2.
\]
The two summands have disjoint odd and even supports. Thus both
vanish exactly under the conditions in \eqref{eq:descent}.
Each tail mirror pair contains one even and one odd position,
which proves the assertion about zeros.
\end{proof}

\begin{lemma}[Two missing coefficients and translation]
\label{lem:two-zero-periodic}
Let $d$ be positive and even, and $1\le a<b\le d-1$. Then
$e=S_d+x^a+x^b$ satisfies $\Delta e=0$ exactly in the following cases:
\begin{equation}\label{eq:two-zero-types}
\begin{array}{c|c|c}
 &d&\{a,b\}\\ \hline
\mathrm{I}&2^r-2,\ r\ge3&\{2^u,2^v\},\quad0\le u<v\le r-1\\
\mathrm{II}&2^r,\ r\ge2&\{2^r-1,2^t\},\quad0\le t\le r-1\\
\mathrm{III}&2^r+2^s-2,\ r>s\ge2&\{2^r-1,2^s-1\}.
\end{array}
\end{equation}
\end{lemma}
\begin{proof}
A binary submask of an integer is obtained by retaining any subset of
its nonzero binary digits; it is proper when it differs from the integer.
The binary expansion identity
$(x+1)^v=\prod_{j:\,v_j=1}(x^{2^j}+1)$ shows that
$\Delta x^v$, for $v>0$, consists of the proper binary submasks
of $v$, and has $2^{\operatorname{wt}_2(v)}-1$ terms.
Consequently $\Delta x^v=1$ exactly for powers of two.
For distinct positive $v,w$, equality $\Delta x^v=\Delta x^w$
also occurs exactly when both are powers of two: with at least
two binary digits, the singleton submasks recover all set digits.
We will also use
\begin{equation}\label{eq:geometric-translation}
 x(x+1)\Delta S_t=x^{t+2}+(x+1)^{t+2}+1.
\end{equation}
It follows that $\Delta S_t=0$ exactly when $t+2$ is a power of two.

For $0\le b\le t$, the one-missing-coefficient polynomial
$S_t+x^b$ is translation invariant exactly when
\[
 (t,b)=(2^s-2,0)\quad\hbox{or}\quad(2^s-1,2^s-1),\qquad s\ge1.
\]
Indeed equality of its values at $0,1$ forces $b=0$ if $t$ is even.
If $t$ is odd, translation invariance forces its leading term to
cancel, since a nonzero invariant polynomial has even degree, so
$b=t$. Equation~\eqref{eq:geometric-translation} completes both cases,
including the zero polynomial at $t=b=0$.

Put $d=2u$ and write $e=E^2+xO^2$.
Odd and even coefficient comparison gives
$\Delta e=0\Longleftrightarrow\Delta O=0,\ \Delta E=O$.
The useful residual is
\[
 B_u=S_{u-1}+\Delta S_u
     =\frac{x^{u+1}+(x+1)^{u+1}+1}{x}.
\]
If $a,b$ are even, then $O=S_{u-1}$. Thus $u+1$ is a power
of two and $B_u=0$. The remaining condition is
$\Delta(x^{a/2}+x^{b/2})=0$, so $a/2,b/2$ are powers of two.
These are precisely the even-even pairs in type I.

If $a,b$ are odd, then $E=S_u$ and
$B_u=x^{(a-1)/2}+x^{(b-1)/2}$. Its two terms force
$\operatorname{wt}_2(u+1)=2$. Writing $u+1=2^v+2^w$, $v>w\ge0$,
gives $B_u=x^{2^v-1}+x^{2^w-1}$. This yields type III when $w>0$;
$w=0$ yields type II with $t=0$.

Finally suppose the missing positions are $2A$ and $2B+1$.
Then $O=S_{u-1}+x^B$ and $B_u=x^B+\Delta x^A$.
The one-missing-coefficient classification gives either
$u=2^s-1,B=0$, or $u=2^s,B=u-1$.
In the first case $B_u=0$; in the second $B_u=x^{u-1}+1$.
Both force $\Delta x^A=1$, so $A$ is a power of two.
They give respectively type I with one missing position $1$,
and type II with $t>0$. These cases exhaust the parities.
All displayed pairs satisfy the required equations on substitution,
so the classification is sufficient as well as necessary.
\end{proof}

\begin{theorem}[Complete four-zero intersection]\label{thm:quadratic}
For the candidates \eqref{eq:candidate},
$\mathcal M_n(f_{n,i,j})=f_{n,i,j}$ if and only if one of the
following holds:
\begin{align}
 n&=2^R+2^S-2,\quad R>S\ge2,&
 (i,j)&=(2^S-2,2^S-1);\label{eq:quadratic-A}\\
 n&=2^R+2,\quad R\ge3,&
 (i,j)&=(3,2^{R-1});\label{eq:quadratic-B}\\
 n&=14,&&
 (i,j)\in\{(2,4),(2,5),(4,5)\}.\label{eq:quadratic-C}
\end{align}
Equivalently, these are exactly the four-zero polynomials in the
image of $\mathcal T$.
\end{theorem}
\begin{proof}
We first use coefficient comparison to reduce the four-zero condition to
the two-zero polynomial in Lemma~\ref{lem:descent}. We then apply the
three types in Lemma~\ref{lem:two-zero-periodic} and impose the second
descent identity on each type.

Write $n=2k$. The coefficient of $x^{n-1}$ in
$\mathcal M_n(f)$ is $f'(1)=k-2$ in $\mathbb F_2$, since the
four missing tail positions include two odd positions: each pair
$\{r,n-1-r\}$ contains one odd position because $n-1$ is odd.
Its value in $f$ is $1$, so $k$ is odd. Comparison of the
$x^2$ coefficients gives $[x]f=\binom{2k}{2}=k=1$.
Thus positions $1,n-2$ are not missing. In Lemma~\ref{lem:descent},
$e=S_d+x^a+x^b$ is therefore monic of even degree $d=k-1$,
with $1\le a<b\le d-1$.

By Lemma~\ref{lem:two-zero-periodic}, $e$ has one of the three types
in \eqref{eq:two-zero-types}. The second descent condition is
equivalent to
\begin{equation}\label{eq:second-descent}
\begin{split}
 x^a+x^b+\Delta x^{d+1-a}+\Delta x^{d+1-b}
 =\frac{x^{d+2}+(x+1)^{d+2}+1}{x}.
\end{split}
\end{equation}
For type I the right side is zero. The binary masks $d+1-a$ and
$d+1-b$ have all $r$ digits set except the distinct digits specified
by $a,b$. Their proper-submask supports have symmetric difference
of size $2^{r-1}-2$: among all submasks the difference has size
$2^{r-1}$, and deleting the two top masks removes two terms.
Equation~\eqref{eq:second-descent} requires that difference to be
the two terms $x^a+x^b$, so $r=3$. For $r=3$ it is exactly that
pair, giving $d=6$ and the three cases \eqref{eq:quadratic-C}.

For type II put $\{a,b\}=\{d-1,2^t\}$. The equation becomes
\[
 \Delta x^{d+1-2^t}=x^{2^t}+x+1.
\]
For $t=0$ both sides are $1$. For $t\ge1$ the right side has
three terms, while $d+1-2^t$ has $r-t+1$ set digits. Thus
$2^{r-t+1}-1=3$, forcing $t=r-1$, which also suffices.
The two choices yield respectively \eqref{eq:quadratic-A} with
$S=2$, and \eqref{eq:quadratic-B}.

For type III the two $\Delta$ terms are both $1$ and cancel.
The right side is $x^a+x^b$, so every such pair satisfies the
condition. Restoring the even zero positions $2a,2b$ and their
tail mirrors gives \eqref{eq:quadratic-A} with $R=r+1,S=s+1\ge3$.
All these steps are reversible, proving the assertion in every degree.
The equivalent image statement follows from the classical
representation preceding Lemma~\ref{lem:descent}.
\end{proof}

\begin{corollary}[Reducibility and failure of iteration]
\label{cor:quadratic-limits}
Every member of \eqref{eq:quadratic-A} is reducible.
If $h$ is itself a four-zero polynomial of the form
\eqref{eq:candidate}, then $\mathcal T(h)$ is never a member of
that family. Among the three polynomials in
\eqref{eq:quadratic-C}, only $f_{14,2,4}$ is irreducible.
\end{corollary}
\begin{proof}
For $f=\mathcal T(h)$ let $m=n/2$; its constant term gives $h(0)=1$,
and the $x^2$ coefficient gives
\begin{equation}\label{eq:quadratic-trace}
 [x]h=[x^2]f+\binom m2.
\end{equation}
In \eqref{eq:quadratic-A}, $S=2$ gives $m\equiv1\pmod4$ and
$[x^2]f=0$, while $S\ge3$ gives $m\equiv3\pmod4$ and $[x^2]f=1$.
Hence $[x]h=0$ in either case. If $h$ is reducible, multiplicativity
of the transform makes $f$ reducible. If $h$ is irreducible, the
classical binary $Q$ criterion \cite[Theorem 1.1]{Ugo13} says that
$h^Q$, and hence $f=h^Q(x+1)$, is irreducible exactly when $[x]h=1$.
This proves the first assertion.

The theorem, or its coefficient proof, gives $n\equiv2\pmod4$
for every four-zero output. A four-zero input has even degree,
so its transformed degree is divisible by four. This proves the
iteration obstruction for this particular transform.

For the exceptional pairs $(2,4),(2,5),(4,5)$ the corresponding
$h$ are, respectively,
\[
 x^7+x^6+x^3+x+1,\quad
 x^7+x^6+x^3+x^2+x+1,\quad
 x^7+x^6+x^4+x^2+1.
\]
The first satisfies $x^{128}\equiv x\pmod h$ and
$\gcd(h,x^2+x)=1$, so is irreducible by the prime-degree Rabin test;
its linear coefficient is $1$. The second has root $1$,
and the third has zero linear coefficient. The same $Q$ criterion
finishes the proof.
\end{proof}

\begin{proposition}[Explicit descent for the B family]\label{prop:B-descent}
For $R\ge3$, put $q=2^{R-1}$. The unique polynomial with
$\mathcal T(h_R)=f_{2q+2,3,q}$ is
\begin{equation}\label{eq:B-descent}
 h_R(x)=1+x+x^2+(1+x)\left(x^q+
       \sum_{t=2}^{R-2}x^{q-2^t+1}\right),
\end{equation}
where the sum is empty for $R=3$. Moreover, $\deg h_R=q+1$ and
$[x]h_R=1$.
\end{proposition}
\begin{proof}
Set $L(x)=\sum_{a=0}^{R-2}x^{2^a}$ and
\[
 F(x)=x^{q+2}h_R(1/x)
     =(x^2+x+1)x^q+(x+1)(L(x)+x^2).
\]
For $s=X^{-1}$ and $w=s+s^2$, telescoping gives $L(w)=s+s^q$;
substitution and collection of terms give
\[
 s^{-2q-4}F(w)=(X+1)f_{2q+2,3,q}(X).
\]
The left side is $(X+1)\mathcal T(h_R)(X)$, proving the formula;
the distinct leading degrees of the transform basis give uniqueness.
The exponents in the sum lie between $q/2+1$ and $q-3$, so the
degree and linear coefficient follow directly.
\end{proof}
The classical $Q$ criterion makes irreducibility of $h_R$ equivalent
to that of its output.

\subsection{The translation branch}\label{sec:translation}
We next impose $f(x+1)=f(x)$. The invariant representation
$f=H(x^2+x)$ and its Artin--Schreier irreducibility criterion are
classical \cite[Lemma 2.1 and Theorem 2.2]{Kap17}. The additional
constraints below come from the tail palindrome and its missing
coefficients.

\begin{lemma}[Degree compatibility]\label{lem:translation-degree}
If a monic tail-palindromic polynomial $f$ of even degree $n\ge6$
satisfies $f(x+1)=f(x)$, then $n+2$ is a power of two.
\end{lemma}
\begin{proof}
Let $F(X,Y)=Y^nf(X/Y)$ and $N=n+1$. Tail symmetry gives
\[
 XF(X,Y)+YF(Y,X)=X^N+Y^N.
\]
By translation invariance, put
$A=F(X,Y)=F(X+Y,Y)$, $B=F(Y,X)=F(X+Y,X)$ and
$C=F(X,X+Y)=F(Y,X+Y)$. Applying the displayed identity to
three pairs of variables yields
\[
\begin{aligned}
 XA+YB&=X^N+Y^N,\\
 (X+Y)A+YC&=(X+Y)^N+Y^N,\\
 XC+(X+Y)B&=X^N+(X+Y)^N.
\end{aligned}
\]
Multiply these equations by $X+Y$, $X$ and $Y$, respectively,
and add. All terms on the left cancel, leaving
$(X+Y)^{n+2}=X^{n+2}+Y^{n+2}$. In characteristic two this
forces $n+2$ to be a power of two: if its binary expansion has
at least two nonzero digits, a singleton proper submask gives
a nonzero mixed binomial term.
\end{proof}

\begin{proposition}[A missing-coefficient obstruction]
\label{prop:translation-zero-obstruction}
Let $f$ be monic of even degree $n=2k\ge6$, with constant term one
and a tail palindrome containing $2s$ zeros. If $f(x+1)=f(x)$,
write $f=H(x^2+x)$ with $H$ monic of degree $k$. Then
\begin{equation}\label{eq:translation-zero-trace}
 [x^{k-1}]H=s\pmod2.
\end{equation}
In particular, if the number of tail zeros is divisible by four,
then $f$ is reducible.
\end{proposition}
\begin{proof}
With $e,o$ as in Lemma~\ref{lem:descent}, tail symmetry gives
$f=(x^k+e)^2+xo^2$. Odd and even coefficient comparison gives
\begin{equation}\label{eq:translation-descent}
 \Delta f=0\quad\Longleftrightarrow\quad
 \Delta o=0,\qquad \Delta(x^k+e)=o.
\end{equation}
Indeed $\Delta f=x(\Delta o)^2+
[\Delta(x^k+e)+o(x+1)]^2$.
The formal coefficient vector of $o$ has length $k$ and $s$ zeros,
so $[x^{n-2}]f=o(0)=o(1)=k-s$ in $\mathbb F_2$.
Lemma~\ref{lem:translation-degree} gives
$k=2^{R-1}-1\equiv3\pmod4$, where $R\ge3$.
The classical coefficient comparison in $H(x^2+x)$ therefore gives
\[
 [x^{k-1}]H=[x^{2k-2}]f+\binom{k}{2}=s\pmod2.
\]
Now suppose that $s$ is even. If $H$ is reducible, so is its composition.
Otherwise let $\beta$ be a root of the irreducible $H$ in
$\mathbb F_{2^k}$. Its absolute trace is $[x^{k-1}]H=0$,
so $z^2+z=\beta$ has a solution in that same field.
Thus $f$ has a root of degree at most $k$ and is reducible.
The last step is classical: specialize \cite[Theorem~2.2(2)]{Ago77}
to $p=2$, $n=r=1$ and $a=1$ in Agou's notation, with irreducible
outer polynomial $H$; see also \cite[Theorem~2.2]{Kap17}.
The same binary specialization follows from Long
\cite[Theorems~4.1--4.2, pp.~69 and 71]{Long73}; the basic
Artin--Schreier trace criterion is classical, as noted in
\cite[p.~63]{Long73}. The associated high-coefficient restriction appears in
\cite[Theorem~2.8(3)]{Kap17}.
\end{proof}

The restriction $n\ge6$ is essential: $x^2+x+1$ is a
translation-invariant irreducible tail palindrome with no zeros.
For odd $s$, \eqref{eq:translation-zero-trace} only removes this
trace obstruction; it does not establish irreducibility of $H$.

\begin{theorem}[Complete translation intersection]
\label{thm:translation-classification}
The polynomial $f_{n,i,j}$ in \eqref{eq:candidate} satisfies
$f_{n,i,j}(x+1)=f_{n,i,j}(x)$ if and only if
\begin{equation}\label{eq:translation-four}
 (n,i,j)\in\{(14,3,4),(30,5,9),(30,5,12),(30,9,12)\}.
\end{equation}
All four polynomials are reducible.
\end{theorem}
\begin{proof}
In \eqref{eq:translation-descent}, put $d=k-1$. The polynomial
$o$ is monic of degree $d$ and has two missing coefficients.
Since $k$ is odd, $o(1)=k-2=1$; invariance gives $o(0)=1$.
Thus $o=S_d+x^a+x^b$ with $1\le a<b\le d-1$, and
$e=S_d+x^{d-a}+x^{d-b}$.
By Lemma~\ref{lem:translation-degree}, $d=2^r-2$, $r\ge2$.
The case $r=2$ has only one interior position and is impossible.
For $r\ge3$, Equation~\eqref{eq:geometric-translation} gives
$\Delta S_d=0$. Hence $\Delta x^a=\Delta x^b$, which, by the
binary-support observation in the proof of
Lemma~\ref{lem:two-zero-periodic}, forces $a,b$ to be distinct
powers of two. The second condition in
\eqref{eq:translation-descent}, using $\Delta x^{d+1}=S_d$, becomes
\begin{equation}\label{eq:translation-support}
 \Delta x^{d-a}+\Delta x^{d-b}=x^a+x^b.
\end{equation}

If $a,b$ are both even, take square roots in this equation.
The two exponents on the left are $(r-1)$-digit all-one binary
masks with different singleton digits removed. Their full
submask supports have symmetric difference of size $2^{r-2}$;
removing the two top masks leaves $2^{r-2}-2$ terms.
The right side has two terms, so $r=4$.
For three-digit masks the remaining difference is precisely the
two removed singleton digits; thus all pairs from $\{2,4,8\}$
satisfy \eqref{eq:translation-support}.

Otherwise $a=1$ and $b$ is an even power of two. The term
$\Delta x^{d-b}$ has only even exponents. The mask
$d-1=2^r-3$ has all digits except digit one set, so
$\Delta x^{d-1}$ has $2^{r-2}-1$ odd terms. The right side has
only the odd term $x$, forcing $r=3$.
Here $\Delta x^5=x^4+x+1$ and $\Delta x^4=\Delta x^2=1$;
of $b=2,4$, exactly $b=4$ satisfies the equation.

A missing position $a$ in $o$ corresponds to the two positions
$2a+1$ and $n-2-2a$ in $f$. Thus $d=6$, $\{a,b\}=\{1,4\}$
gives $(14,3,4)$; $d=14$ and the three pairs from $\{2,4,8\}$
give the degree-$30$ cases in \eqref{eq:translation-four}.
The substitutions above also prove sufficiency.
Reducibility follows uniformly from
Proposition~\ref{prop:translation-zero-obstruction} with $s=2$.
\end{proof}

\subsection{Full-degree quadratic constructions}\label{sec:quadratic-scope}
Here a full-degree quadratic pullback over $\mathbb F_2$ means
\[
 F=h^mP(g/h),\qquad \gcd(g,h)=1,\qquad
 \max(\deg g,\deg h)=2,
\]
where $\deg P=m>1$ and $\deg F=2m$.
The last condition is automatic for irreducible
$P$; in general the degree drops exactly when $h_2\ne0$ and
$P(g_2/h_2)=0$ \cite[Section 2]{MP17}.

\begin{corollary}[Boundary for full-degree quadratic pullbacks]
\label{cor:all-quadratic}
If a full-degree quadratic pullback is a four-zero tail palindrome,
then it belongs either to the lists in Theorem~\ref{thm:quadratic}
or to \eqref{eq:translation-four}, and its degree is $2$ modulo $4$.
If it is irreducible, it is either $f_{14,2,4}$ or a member of
\eqref{eq:quadratic-B}. Consequently, no sequence of quadratic
pullbacks, each doubling the degree, can reach a four-zero tail
palindrome after starting at an even degree.
\end{corollary}
\begin{proof}
Write $g=\sum_{j=0}^2g_jx^j$, $h=\sum_{j=0}^2h_jx^j$, and put
\[
 a=g_2h_1+g_1h_2,\quad b=g_0h_2+g_2h_0,\quad
 c=g_1h_0+g_0h_1.
\]
The classical quadratic framework gives $b^2+ac=1$ and, in the
separable case,
\[
 (ax+b)^{2m}F\!\left(\frac{bx+c}{ax+b}\right)=F(x)
 \qquad\text{\cite[Section 4, Lemma 9]{MP17}}.
\]
In the separable case the possible triples are exactly $(1,0,1),(0,1,1),(1,1,0)$,
corresponding to reciprocity, translation and $\mathcal M_{2m}$.
There is no nontrivial scalar normalization over $\mathbb F_2$.
In the inseparable case $g_1=h_1=0$, the output is a square;
its odd coefficients vanish, whereas a four-zero candidate has
$[x^{2m-1}]F=1$.

The reciprocal branch contains no four-zero candidate: combining
$f^*=f$ with Lemma~\ref{lem:reciprocal-multiplicity} gives
$(x+1)f=x^{n+1}+1=(x+1)S_n$, hence $f=S_n$.
The other two branches are classified by
Theorems~\ref{thm:quadratic} and~\ref{thm:translation-classification}.
The irreducible alternatives follow from their reducibility
conclusions. All listed degrees are $2$ modulo $4$.
Starting at an even degree, every subsequent doubled degree is
divisible by four, proving the last assertion.
\end{proof}

The three involutions thus connect the support symmetries to an
exhaustive list of full-degree quadratic outputs.


\section{Primitive counts and reproducible computations}\label{sec:computations}
Complete enumeration of the $161700$ candidates at the $98$ even degrees
$6\le n\le200$ gives
\[
 \sum_{\substack{6\le n\le200\\ n\text{ even}}}N_4(n)=2663,
 \qquad N_4(n)>0\ \text{at every tested degree}.
\]
A separate implementation reproduces every complete count and first pair.
At the $50$ even degrees $202\le n\le300$, checking the lexicographic
prefixes covers $5410$ further candidates and verifies the first pair
at each degree.

\begin{table}[ht]
\centering\small
\begin{tabular}{rrrr}
\toprule
$n$ & candidates & $N_4(n)$ & first $(i,j)$\\
\midrule
6 & 1 & 1 & $(1,2)$\\
8 & 3 & 1 & $(2,3)$\\
10 & 6 & 2 & $(1,2)$\\
12 & 10 & 3 & $(1,3)$\\
20 & 36 & 1 & $(4,5)$\\
40 & 171 & 11 & $(1,4)$\\
62 & 435 & 16 & $(1,17)$\\
100 & 1176 & 31 & $(1,23)$\\
126 & 1891 & 55 & $(1,2)$\\
154 & 2850 & 45 & $(1,23)$\\
194 & 4560 & 20 & $(6,65)$\\
200 & 4851 & 17 & $(1,71)$\\
\bottomrule
\end{tabular}
\caption{Selected complete counts. All $98$ rows and all successful
position pairs through degree $200$ are supplied with the verification data.}
\label{tab:counts}
\end{table}

\paragraph{Arithmetic verification and reproducibility.}
The verification code constructs coefficients directly and tests
irreducibility, followed by the standard exact-order conditions
\begin{equation}\label{eq:ordertest}
 x^{2^n-1}\equiv1\pmod f,\qquad
 x^{(2^n-1)/p}\not\equiv1\pmod f\quad(p\mid2^n-1\text{ prime}).
\end{equation}
The arithmetic tests use python-flint $0.9.0$ with FLINT $3.6.0$
\cite{Flint}.
The $148$ order factorizations correspond to the $98$ complete-count
degrees and the $50$ further first-witness degrees. Exact division leaves
residual cofactor one; all $484$ distinct prime factors, of at most $144$
bits, pass FLINT's rigorous \texttt{fmpz\_is\_prime} procedure.
A standard-library implementation independently agrees on all $220$
candidates through degree $24$.

The folder \path{verification/primitivity/} supplies complete candidate
verdicts, successful positions, exact order-test residues for the $148$
first witnesses, recovered integer factorizations and input hashes.
Historical inputs and correction records are preserved with the audit.
The data and verification programs are publicly archived at
\href{https://doi.org/10.5281/zenodo.23766722}{doi:10.5281/zenodo.23766722}
and in the
\href{https://github.com/63331849121/yijing-mathematical-structures/tree/paperF-v2026.10.29-r1/PaperF_fourzero}{fixed GitHub release}.

The companion programs also check the reciprocal filter, both quadratic
classifications and the descent identities; detailed ranges accompany
the archived outputs.
The four translation outputs have independently certified factor degrees
$2+12$, $12+18$, $6+24$ and $7+7+8+8$, with a repeated factor in the last.
The B-sequence factorization data cover $3\le R\le10$:
the first term is irreducible and the next seven reducible.

\section{Structural consequences and open questions}\label{sec:limits}
The reciprocal screen counts positions compatible with reciprocal
coprimality. The quadratic classifications determine every full-degree
output and yield a uniform obstruction from translation symmetry.

\paragraph{Irreducibility in the remaining branch.}
Within full-degree quadratic constructions, the B family
\eqref{eq:quadratic-B} is the remaining irreducibility question in
Corollary~\ref{cor:all-quadratic}. Which polynomials $h_R$ of
Proposition~\ref{prop:B-descent} are irreducible? The explicit descent and the
binary $Q$ criterion make this equivalent to classifying the irreducible
four-zero outputs in that branch.

\paragraph{Primitive existence and further constructions.}
Conjecture~\ref{conj:C4} asks for primitive members at every even degree,
linking the prescribed support to the multiplicative order of roots.
The complete counts and first witnesses give concrete instances; the
quotient and finite-field descriptions in
Appendix~\ref{sec:classical-descriptions} express the support condition
and the order condition explicitly. Higher-degree maps, selected factors and degree-dropping
substitutions provide further construction questions.

\appendix
\section{Coefficient and finite-field descriptions}\label{sec:classical-descriptions}
The coefficient quotient organizes position pairs into classes with explicit
multiplicities; the finite-field parametrization encodes primitivity
through coprime order components.

\subsection{Coefficient quotient}\label{sec:quotient}
Let $\rho$ reverse $\mathbb F_2^k$, and set
\[
 P=\{(A,\rho A):A\in\mathbb F_2^k\},\quad
 K=\{(C,C):C=\rho C\},\quad U=P/K,\quad r=\lfloor k/2\rfloor.
\]
The map $[(A,\rho A)]\mapsto A+\rho A$ identifies $U$ with
$\operatorname{im}(I+\rho)$, of dimension $r$.
Prepending a zero to $x\in\mathbb F_2^{k-1}$ defines
$\Phi(x)=[(\widetilde x,\rho\widetilde x)]$.
The all-ones tail lies in $K$, so deleting positions $i,j$ gives
class $\Phi(e_i+e_j)$. With zero-based coordinates and a suitable basis
$z_1,\ldots,z_r$, the columns $u_i=\Phi(e_i)$ form the multiset
\begin{equation}\label{eq:columns}
 \{u_1,\ldots,u_{k-1}\}=
 \begin{cases}
 \{z_1,z_1,\ldots,z_{r-1},z_{r-1},z_r\},&k=2r,\\
 \{z_1,z_1,\ldots,z_{r-1},z_{r-1},z_r,0\},&k=2r+1.
 \end{cases}
\end{equation}
Indeed pair positions $i,k-1-i$; the boundary pair contributes the
singleton $z_r$ and an unpaired middle position contributes zero.
In particular $\Phi$ is onto.

For $s=\sum s_a z_a$, write $h=\sum_{a<r}s_a$ as an integer and
$\varepsilon=k\bmod2$. The ordinary coset weight enumerator is
\begin{equation}\label{eq:enumerator}
 A_s(z)=\sum_{\Phi(x)=s}z^{\operatorname{wt}(x)}
       =(1+z)^\varepsilon z^{s_r}(2z)^h(1+z^2)^{r-1-h}.
\end{equation}
Each repeated column allows $00,11$ for zero syndrome and $01,10$
for nonzero syndrome; the singleton and optional zero column give the
remaining factors. Thus $\mu(s)=[z^2]A_s(z)$ counts four-zero candidates
in class $s$, and the number of occupied classes is
\begin{equation}\label{eq:Dk}
 D(3)=1,\qquad D(k)=1+\binom{\lfloor(k+1)/2\rfloor}{2}\quad(k\ge4).
\end{equation}
This follows by counting the nonzero quadratic coefficients in
\eqref{eq:enumerator}; at $k=3$ the zero class is unoccupied.
For a character $\chi$ of $U$, put
$t=\#\{i:\chi(u_i)=-1\}$ and $\ell=k-1$. Then
\begin{equation}\label{eq:krawtchouk}
 \widehat\mu(\chi)=\sum_{i<j}\chi(u_i)\chi(u_j)
   =K_2(t;\ell)=\frac{(\ell-2t)^2-\ell}{2}.
\end{equation}
These are classical direct-product weight-enumerator and Krawtchouk
identities \cite[Sections 3--4]{Forney11}\cite[Sections 1.1 and 2.1]{Brouwer18}.
The different factorization behavior within a class is visible already
at degree eight:
$x^8+x^7+x^6+x+1$ is primitive, while
$x^8+x^7+x^5+x^2+1=(x^2+x+1)(x^6+x^4+x^2+x+1)$ is reducible;
their tail difference $0110\,0110$ belongs to $K$.

\subsection{Finite-field parametrization}\label{sec:dickson}
Put $q=2^k$, $\mathbb K=\mathbb F_q$, $\mathbb L=\mathbb F_{q^2}$ and
$T=\{\gamma\in\mathbb L^*:\gamma^{q+1}=1\}$.
The coprime orders $q-1,q+1$ give $\mathbb L^*=\mathbb K^*\times T$.
A primitive element is uniquely $\beta=\nu\gamma$ with
$\operatorname{ord}(\nu)=q-1$ and $\operatorname{ord}(\gamma)=q+1$.
Writing $\tau=\gamma+\gamma^{-1}\in\mathbb K$, the standard Dickson
identity $D_a(u+u^{-1})=u^a+u^{-a}$ and Frobenius pairing give
\begin{equation}\label{eq:dickson}
 F_{\nu,\gamma}(x)=\prod_{i=0}^{k-1}
 \left(x^2+\nu^{2^i}D_{2^i}(\tau)x+\nu^{2^{i+1}}\right).
\end{equation}
Here $D_0=0,D_1=t,D_{a+1}=tD_a+D_{a-1}$ in characteristic two
\cite{LN}. Conversely, the two order conditions make this product the
minimal polynomial of a primitive element. Expanding the product and
comparing with \eqref{eq:candidate} gives an exact coefficient-and-order
system in which the support and the two exact-order conditions are
imposed simultaneously.

\begin{thebibliography}{99}
\small
\setlength{\itemsep}{1.5pt}
\setlength{\parsep}{0pt}
\setlength{\parskip}{0pt}
\bibitem{Coh06}
S. D. Cohen,
\emph{Primitive polynomials with a prescribed coefficient},
Finite Fields Appl. \textbf{12} (2006), 425--491.

\bibitem{LN}
R. Lidl and H. Niederreiter,
\emph{Finite Fields},
Encyclopedia of Mathematics and its Applications, vol.~20, 2nd ed., Cambridge University Press, 1997.

\bibitem{Swan62}
R. G. Swan,
\emph{Factorization of polynomials over finite fields},
Pacific J. Math. \textbf{12} (1962), 1099--1106.

\bibitem{FHS86}
H. M. Fredricksen, A. W. Hales, and M. M. Sweet,
\emph{A generalization of Swan's theorem},
Math. Comp. \textbf{46} (1986), 321--331.

\bibitem{Def20}
C. Defant,
\emph{Enumerating cliques in direct product graphs},
J. Combin. \textbf{11} (2020), 351--358.

\bibitem{Brouwer18}
A. E. Brouwer, S. M. Cioab\u{a}, F. Ihringer, and M. McGinnis,
\emph{The smallest eigenvalues of Hamming graphs, Johnson graphs and
other distance-regular graphs with classical parameters},
J. Combin. Theory Ser. B \textbf{133} (2018), 88--121.
\href{https://doi.org/10.1016/j.jctb.2018.04.005}{doi:10.1016/j.jctb.2018.04.005}.

\bibitem{Forney11}
G. D. Forney, Jr.,
\emph{Codes on graphs: duality and MacWilliams identities},
IEEE Trans. Inform. Theory \textbf{57} (2011), 1382--1397.
\href{https://doi.org/10.1109/TIT.2011.2104994}{doi:10.1109/TIT.2011.2104994}.

\bibitem{MP17}
S. Mattarei and M. Pizzato,
\emph{Generalizations of self-reciprocal polynomials},
Finite Fields Appl. \textbf{48} (2017), 271--288.
\href{https://doi.org/10.1016/j.ffa.2017.08.004}{doi:10.1016/j.ffa.2017.08.004}.

\bibitem{Reis20}
L. Reis,
\emph{M\"obius-like maps on irreducible polynomials and rational transformations},
J. Pure Appl. Algebra \textbf{224} (2020), 169--180.
\href{https://doi.org/10.1016/j.jpaa.2019.05.001}{doi:10.1016/j.jpaa.2019.05.001}.

\bibitem{Long73}
A. F. Long, Jr.,
\emph{Factorization of irreducible polynomials over a finite field with
the substitution $x^{p^r}-x$ for $x$},
Duke Math. J. \textbf{40} (1973), 63--76.
\href{https://doi.org/10.1215/S0012-7094-73-04007-6}{doi:10.1215/S0012-7094-73-04007-6}.

\bibitem{Ago77}
S. Agou,
\emph{Factorisation sur un corps fini $\mathbb F_{p^n}$ des
polyn\^omes compos\'es $f(X^{p^r}-aX)$ lorsque $f(X)$ est un
polyn\^ome irr\'eductible de $\mathbb F_{p^n}(X)$},
J. Number Theory \textbf{9} (1977), 229--239.
\href{https://www.sciencedirect.com/science/article/pii/0022314X77900270}{Publisher page}.

\bibitem{Kap17}
G. Kapetanakis,
\emph{Prescribing coefficients of invariant irreducible polynomials},
J. Number Theory \textbf{180} (2017), 615--628.
\href{https://doi.org/10.1016/j.jnt.2017.05.003}{doi:10.1016/j.jnt.2017.05.003}.

\bibitem{Ugo13}
S. Ugolini,
\emph{Sequences of binary irreducible polynomials},
Discrete Math. \textbf{313} (2013), 2656--2662.
\href{https://doi.org/10.1016/j.disc.2013.08.011}{doi:10.1016/j.disc.2013.08.011}.

\bibitem{Flint}
The python-flint and FLINT developers,
\emph{python-flint}, version 0.9.0, with FLINT version 3.6.0.
\href{https://python-flint.readthedocs.io/}{Documentation};
\href{https://flintlib.org/doc/fmpz.html}{integer primality-testing documentation}.
Accessed 27 October 2026.

\end{thebibliography}
\end{document}
